% GENERATED FILE. Edit canonical lessons, then run npm run generate:books.
% canonical-source-hash: fnv1a64:6f7a70deeef6dde8
% canonical-lessons: TE-C179-kammani, TE-C179-cappani, TE-C179-taja, TE-C179-kharidaina, TE-C179-cauka

\chapter{Delicious, Bland, Fresh, Expensive, Cheap}
\label{ch:te-delicious-bland-fresh-expensive-cheap}
\hlchaptermodality{179}

\begin{chapteropening}
\textbf{By the end of this chapter:} \emph{I can say delicious, bland, fresh, expensive and cheap.}

Complete the last lesson of chapter 179: I can say delicious, bland, fresh, expensive and cheap.
\end{chapteropening}

\section[\texorpdfstring{\te{కమ్మని} (kammani)}{కమ్మని (kammani)}]{\te{కమ్మని} (kammani) — delicious, tasty}

\label{lesson:TE-C179-kammani}



\begin{quote}
Before the new one: say the Telugu for a condition, then the Telugu for an excuse.
\end{quote}

\subsection*{You'll want to know: \te{కమ్మని}}
\textbf{\te{కమ్మని}} — \emph{kammani} — \textquotedblleft{}delicious, tasty\textquotedblright{}.

\begin{grammarlens}[title={What you've built}]
One more word for this chapter.
\end{grammarlens}

\subsection*{Your turn}
\begin{itemize}
\raggedright
  \item \emph{Say it:} \emph{kammani}
  \item \emph{Say it:} \emph{kammani}, once more
  \item \emph{Say it:} say \emph{kammani}
  \item \emph{Recall:} say the Telugu for a condition, then the Telugu for an excuse, then say \emph{kammani} again
\end{itemize}

\begin{tcolorbox}[breakable,colback=teal!4,colframe=teal!35!black,title={Before you move on}]
What does \te{కమ్మని} mean? (\textquotedblleft{}Delicious, tasty\textquotedblright{}.) Say it once more.
\end{tcolorbox}
\section[\texorpdfstring{\te{చప్పని} (cappani)}{చప్పని (cappani)}]{\te{చప్పని} (cappani) — bland, tasteless}

\label{lesson:TE-C179-cappani}



\begin{quote}
Before the new one: say the Telugu for an excuse, then the Telugu for delicious.
\end{quote}

\subsection*{You'll want to know: \te{చప్పని}}
\textbf{\te{చప్పని}} — \emph{cappani} — \textquotedblleft{}bland, tasteless\textquotedblright{}.

\begin{grammarlens}[title={What you've built}]
One more word for this chapter.
\end{grammarlens}

\subsection*{Your turn}
\begin{itemize}
\raggedright
  \item \emph{Say it:} \emph{cappani}
  \item \emph{Say it:} \emph{cappani}, once more
  \item \emph{Say it:} say \emph{cappani}
  \item \emph{Recall:} say the Telugu for an excuse, then the Telugu for delicious, then say \emph{cappani} again
\end{itemize}

\begin{tcolorbox}[breakable,colback=teal!4,colframe=teal!35!black,title={Before you move on}]
What does \te{చప్పని} mean? (\textquotedblleft{}Bland, tasteless\textquotedblright{}.) Say it once more.
\end{tcolorbox}
\section[\texorpdfstring{\te{తాజా} (tājā)}{తాజా (tājā)}]{\te{తాజా} (tājā) — fresh}

\label{lesson:TE-C179-taja}



\begin{quote}
Before the new one: say the Telugu for delicious, then the Telugu for bland.
\end{quote}

\subsection*{You'll want to know: \te{తాజా}}
\textbf{\te{తాజా}} — \emph{tājā} — \textquotedblleft{}fresh\textquotedblright{}.

\begin{grammarlens}[title={What you've built}]
One more word for this chapter.
\end{grammarlens}

\subsection*{Your turn}
\begin{itemize}
\raggedright
  \item \emph{Say it:} \emph{tājā}
  \item \emph{Say it:} \emph{tājā}, once more
  \item \emph{Say it:} say \emph{tājā}
  \item \emph{Recall:} say the Telugu for delicious, then the Telugu for bland, then say \emph{tājā} again
\end{itemize}

\begin{tcolorbox}[breakable,colback=teal!4,colframe=teal!35!black,title={Before you move on}]
What does \te{తాజా} mean? (\textquotedblleft{}Fresh\textquotedblright{}.) Say it once more.
\end{tcolorbox}
\section[\texorpdfstring{\te{ఖరీదైన} (kharīdaina)}{ఖరీదైన (kharīdaina)}]{\te{ఖరీదైన} (kharīdaina) — expensive}

\label{lesson:TE-C179-kharidaina}



\begin{quote}
Before the new one: say the Telugu for bland, then the Telugu for fresh.
\end{quote}

\subsection*{You'll want to know: \te{ఖరీదైన}}
\textbf{\te{ఖరీదైన}} — \emph{kharīdaina} — \textquotedblleft{}expensive\textquotedblright{}.

\begin{grammarlens}[title={What you've built}]
One more word for this chapter.
\end{grammarlens}

\subsection*{Your turn}
\begin{itemize}
\raggedright
  \item \emph{Say it:} \emph{kharīdaina}
  \item \emph{Say it:} \emph{kharīdaina}, once more
  \item \emph{Say it:} say \emph{kharīdaina}
  \item \emph{Recall:} say the Telugu for bland, then the Telugu for fresh, then say \emph{kharīdaina} again
\end{itemize}

\begin{tcolorbox}[breakable,colback=teal!4,colframe=teal!35!black,title={Before you move on}]
What does \te{ఖరీదైన} mean? (\textquotedblleft{}Expensive\textquotedblright{}.) Say it once more.
\end{tcolorbox}
\section[\texorpdfstring{\te{చౌక} (cauka)}{చౌక (cauka)}]{\te{చౌక} (cauka) — cheap}

\label{lesson:TE-C179-cauka}



\begin{quote}
Before the new one: say the Telugu for fresh, then the Telugu for expensive.
\end{quote}

\subsection*{You'll want to know: \te{చౌక}}
\textbf{\te{చౌక}} — \emph{cauka} — \textquotedblleft{}cheap\textquotedblright{}.

\begin{grammarlens}[title={What you've built}]
One more word for this chapter.
\end{grammarlens}

\subsection*{Your turn}
\begin{itemize}
\raggedright
  \item \emph{Say it:} \emph{cauka}
  \item \emph{Say it:} \emph{cauka}, once more
  \item \emph{Say it:} say \emph{cauka}
  \item \emph{Recall:} say the Telugu for fresh, then the Telugu for expensive, then say \emph{cauka} again
\end{itemize}

\begin{tcolorbox}[breakable,colback=teal!4,colframe=teal!35!black,title={Before you move on}]
What does \te{చౌక} mean? (\textquotedblleft{}Cheap\textquotedblright{}.) Say it once more.
\end{tcolorbox}
